\begin{qfdata}[Trading reference data: universe, symbology, corporate actions, static and exchange data, condition codes, special days, futures and options reference, event calendars, related and composite instruments]\label{data:trading-reference-data}
\qffield{Statement} Velu, Hardy and Nehren call reliable reference data ``the key foundation of a robust quantitative strategy development'' and list what a research platform must hold \citep[chap.~1, sec.~1.4.1, pp.~26--35]{s3-velu}:
\begin{itemize}\setlength{\itemsep}{0pt}
\item the \emph{trading universe}, which changes daily with listings and delistings; knowing when a stock stopped trading avoids survivor bias (p.~26);
\item \emph{symbology mapping} (ISIN, SEDOL, RIC, Bloomberg tickers) and \emph{ticker changes}: the mapping ``needs to persist as point-in-time data and allow for historical `as of date' usage, requiring the implementation of a bi-temporal data structure'', handling changes backward and forward (pp.~26--27);
\item \emph{corporate actions}: dividends (announcement and payment dates), splits, reverse splits, rights offers, mergers, spin-offs, changes in free float or shares outstanding, suspensions; prices and volumes must be adjusted backward for splits, and suspensions leave gaps that affect backtests (pp.~27--28);
\item \emph{static data}: country, sector (GICS, ICB, TOPIX's 33 sectors), primary exchange, quotation currency and quote factor, with USD price $=$ local price $\times$ FX $\times$ quote factor and cent-quoted currencies such as GBp (pp.~28--29);
\item \emph{exchange data}: holiday calendars, session hours with auctions, lunch breaks and daylight-saving shifts, disrupted days, tick sizes, trade and quote lots, limit-up and limit-down rules and short-sale restrictions, kept historically ``in order to run realistic historical backtests'' (pp.~29--30);
\item \emph{condition codes} on trades and quotes, mapped so that, for example, off-exchange trades reported for regulatory purposes are excluded from accessible liquidity and from daily volume, high and low (p.~31);
\item \emph{special-day calendars}: half days, local holidays and exam days, month and quarter ends, index rebalance days with volume skewed to the close, and expiry days (pp.~31--32);
\item \emph{futures and options reference data}: expiry calendars, the most liquid contract, roll dates, contract size and quotation, session phases, and option chains (pp.~32--34);
\item \emph{event calendars} for macro releases and company events, \emph{related tickers} (composite tickers, dual listings, ADRs and GDRs), \emph{composite assets} with their constituents, cash components and divisors, and \emph{latency tables} between data centres (pp.~34--35).
\end{itemize}

\qffield{Conditions} Examples as of about 2018--2020; exchange rules and calendars change. Every table that changes over time must be stored with the dates on which each value applied and was known, so that a backtest sees the universe, identifiers, corporate actions, tick sizes and rules in force at each decision time. Velu, Hardy and Nehren's reasons for special days are practical: a single special day moves a 30- or 60-day ADV, and low-order autoregressive models underperform on and after such days, so modellers ``often remove special days and model normal days first'' and model special days separately (p.~32).

Futures liquidity moves between contracts at the roll (about 1.5 million E-mini contracts a day in the front month against 30,000 in the next before expiry, reversing in the last days, p.~33), so rolling-window statistics must blend contracts across roll dates.

\qffield{Practice} Velu, Hardy and Nehren's ABC$\to$DEF example (p.~27) shows the three queries a symbology table must answer: a 30-day history as of a date before the change (old ticker), after it (both tickers spliced), and a forward window for a position held across it. Examples they give: European equity volume 30--40 percent lower on U.S. holidays; Indonesia's longer Friday lunch break; U.S. daylight saving starting about two weeks before Europe's (pp.~29--30). These are illustrations as of the book's writing, not current calendars.

\qffield{Pitfalls} Velu, Hardy and Nehren distinguish the drop in price when the dividend is paid from the volatility that may accompany the announcement (p.~27); interpretation, the price drop occurs on the ex-dividend date. Merger and spin-off data can support event-driven strategies (p.~28); exchange-rule histories must also be retained. The book describes what to collect; it does not test how much results change when it is ignored. A ticker is not an instrument key: tickers are display mnemonics, and a key needs venue, currency, instrument type, expiry, strike and option right where they apply, with corporate-action history and identifier validity dates. Validity dates are the point-in-time requirement above; the attribute list is an interpretation (compare \nolinkurl{exp-gleif} and \nolinkurl{s8-bitemporal}).

\qffield{See also} \hyperref[data:market-data-levels]{market data levels}; \hyperref[data:derived-market-statistics]{derived market statistics}; \hyperref[def:financial-data-types]{financial data types}; \hyperref[meth:etf-trick]{ETF trick} (rolled futures series); \hyperref[ev:index-reconstitution-impact]{index reconstitution impact}; \hyperref[mech:securities-lending-short-sale]{securities lending and short sales}; \hyperref[conv:cme-globex-trading]{CME Globex trading}.

\qffield{Further reading} See \hyperref[meth:bitemporal-reference-tables]{bitemporal reference tables}. On delisting bias, Shumway finds that delistings ``for bankruptcy and other negative reasons are generally surprises'', that ``correct delisting returns are not available for most of the stocks that have been delisted for negative reasons since 1962'', and, from over-the-counter prices, that ``the omitted delisting returns are large'' \citep[abstract]{wide-shumway-delisting}. A universe that keeps delisted stocks but lacks their final returns is therefore still biased upward (Interpretation). Shumway and Warther's abstract, from the publisher's record as served by RePEc, quantifies it for Nasdaq. The missing returns of delisted stocks are ``large and negative on average''. A corrected return of $-55$ percent for missing performance-related delisting returns corrects the bias.

After the correction there is ``no evidence that there ever was a size effect on Nasdaq'' \citep[abstract]{shumway-warther-1999}. Interpretation: a substitute delisting return is an estimate for one exchange and period, to be reported as an assumption. Shumway's full text remains refused. See also CRSP files, Elton, Gruber and Blake on fund survivorship, and Ljungqvist, Malloy and Marston on vendor history rewriting.

\end{qfdata}

\begin{qfdata}[Market data levels: Level I trades and quotes, Level II depth, Level III order messages, and the special order types that complicate them]\label{data:market-data-levels}
\qffield{Statement} Market data are sold ``in various levels of granularity (and price!)'', from exchanges as direct feeds or from vendors \citep[chap.~1, sec.~1.4.2, p.~35]{s3-velu}. \emph{Level I} holds trades and top-of-book quotes, enough to rebuild the best bid and offer; quotes update roughly ten times as often as trades in liquid markets, and trade status and qualifiers (cancelled, late, odd lot, auction, intermarket sweep, average-price, venue) decide whether a print should update the last price and cumulative volume (pp.~35--36).

\emph{Level II} adds every lit order-book update at every level and venue; many users keep five or ten levels (p.~36). \emph{Level III}, message data, gives each order an identifier tracked through execution, cancellation or amendment; once messages from all exchanges are ``consolidated into one single dataset ordered by timestamp'', a full book can be built at any moment (pp.~36--37). Their example of a U.S. Level~III record has timestamp (milliseconds after midnight), ticker, order identifier, message type (add buy, add sell, partial execution, partial cancel, full execution, delete, cross, non-displayed execution), shares, price (the last four digits decimals, so ``divide by 10000''), participant identifier and market centre (Table~1.10, p.~37).

\qffield{Conditions} U.S. feed examples as of about 2018--2020; TAQ codes as of the November 2022 specification. Exchange timestamps are matching-engine times; databases also store a reception time, and the latency between them must be accounted for when opportunities are short-lived (p.~36). Special order types change identifiers and priority in venue-specific ways: price-sliding orders ranked at the locking price as hidden and displayed one tick away; pegged and midpoint-peg orders; reserve orders whose displayed part is replenished with a new timestamp; discretionary orders; intermarket sweep orders (pp.~37--38). Velu, Hardy and Nehren advise weighing the value of raw data against simpler binned data, and note that most practitioners outside high-frequency trading build strategies on binned data (pp.~36 and~39).

\emph{Timestamps, sequencing and availability}. Hasbrouck: when data from several markets are combined, ``The quality of the time stamps depends on the accuracy of internal systems clocks, latencies in data processing and transmission, and the nominal reporting precision'', and with several events in one reported second ``the sequencing of reporting is suspect''; for many purposes what matters is ``when knowledge of an event entered the public awareness'', so the timing and sequencing of the real-time consolidated feed are the appropriate ones \citep[chap.~10, p.~105]{s6-hasbrouck}. Time aggregation makes events separated in time look simultaneous (p.~104).

For a feature in a backtest this gives three distinct times: when the event happened at the venue (the matching-engine timestamp), when it reached the consolidated feed, and when the strategy's own systems received it; a feature is available at a decision time only if it was received by then (Interpretation, consistent with the reception-time point above). Corrections and cancellations arrive later as separate records with their own status and qualifiers (above, pp.~35--36), so a feature computed from the corrected history uses information that was not available at the time.

\emph{How TAQ marks them}. NYSE's Daily TAQ trades file, sorted by symbol, time and sequence number, carries a two-digit Trade Correction Indicator: 00 for a regular trade ``not corrected, changed or signified as cancel or error''; 01 for an ``Original trade which was late corrected'', a record that ``contains the original time---HHMM and the corrected data''; 07 and 08 for originals later marked erroneous or cancelled; and 10, 11 and 12 for the cancel, error and correction records that follow them, the last containing ``the correction time and the original `incorrect' data''. Late and out-of-sequence reports are flagged separately in the sale conditions (CTA and UTP codes L ``Sold Last'' and Z ``Sold (Out of Sequence)'', with T and U for extended hours) \citep[v4.0 of 2 November 2022, sec.~3, pp.~17--19]{nyse-daily-taq-spec}.

\emph{Replaying corrections}. In the trades file the \emph{Time} field is when ``the trade was published by the SIP'' (field~1), and a \texttt{12} record carries the correction time \citep[sec.~3, PDF pp.~17--19]{nyse-daily-taq-spec}. A \texttt{01} record, however, places corrected data at the original time. A decision-time panel should therefore take the trade as first reported, with the original values recoverable from the \texttt{12} record, until the correction time. It should keep \texttt{07} and \texttt{08} originals until their error or cancel was published, and treat L and Z prints as known at publication, not at the trade time they report.

\emph{Unresolved times and older editions}. Which time the \texttt{10} and \texttt{11} records carry is not stated. Field~8 defines them only as following the \texttt{07} and \texttt{08} originals, and three earlier editions use the same wording: the specification for 2~October 2006 to 31~July 2012 (V1.0c), v3.0a and v3.1 \citep[V1.0c, PDF p.~19; v3.0a, PDF p.~19; v3.1, PDF p.~17]{nyse-daily-taq-spec-older}. None of the four editions examined states it, so that question stays open. The editions differ elsewhere: V1.0c describes both \texttt{07} and \texttt{08} as originals ``later signified as error'', where later editions have \texttt{08} as cancelled, and v3.0a and v3.1 add ``The final correction will be published''.

\emph{Holden and Jacobsen's screens}. Their sample is 100 stocks drawn at random, 20 from each quintile of trades per day, from April to June 2008 (34 million trades and 351 million quotes). The main-text screens keep quotes from 9:00 a.m.\ and trades from 9:30 a.m.\ to 4:00 p.m. They drop a venue's own crossed or locked quote, a side more than \$2.50 from the previous midpoint when the spread exceeds \$5, quotes with non-normal conditions, and zero or missing prices or sizes, which with the withdrawn-quote adjustment are read as withdrawn sides \citep[published version, sec.~III, p.~1758 and n.~16]{wide-holden-jacobsen}.

NBBO quotes and trades while the NBBO is crossed or locked are deleted (p.~1750). Their DTAQ quote file's ``Quote Cancel/Correction'' field was ``always blank in our sample'' (n.~12, p.~1754). No trade-correction filter is stated in the main text: the paper points to the Internet Appendix for detail (p.~1758).

That appendix prints the July 2013 SAS program for their \emph{Monthly} TAQ recommendations, whose only trade filter keeps ``trades in which the Correction field contains `00' and the Price field contains a value greater than zero'' (\texttt{where corr eq 0 and price gt 0}) \citep[Internet Appendix, sec.~I and the program's step three, PDF pp.~1--2 and~19]{wide-holden-jacobsen}: corrected and cancelled trades are dropped, not replayed. A filter on the final correction code uses after-the-fact information, appropriate for measuring liquidity but not for a decision-time replay (Interpretation). Only this printed program is covered here; their DTAQ code, linked from Holden's page, was not examined.

Their Interpolated Time assigns the $i$th of $I$ events in recorded second $s$ the time $s+(2i-1)/(2I)$ (eq.~(16), p.~1766), and they note that, like matching within the same second, it ``uses future information, extending as much as one second into the future'' (n.~29, p.~1767): an evaluation device for liquidity measures, not a point-in-time rule. With all three adjustments the monthly-TAQ gap to DTAQ vanished for percent quoted spread and shrank by 91 and 97 percent for percent effective spread and price impact (pp.~1750--1751); Jurkatis found the interpolation lowered trade-classification accuracy (\hyperref[meth:trade-sign-classification]{trade-sign classification}).

\qffield{Practice} With Level~III data one can model inter-arrival times, arrival and cancellation rates by distance from the touch, queue position and imbalance, and then ask about the impact of market orders, a limit order's chance of moving up the queue, the probability of earning the spread and the short-horizon direction of prices (p.~38). Interpretation: the levels above usually come with a separate trades stream (price, size, time, venue and condition codes), the industry's L1/L2/L3 labels are Levels~I--III, and a sound rule is to choose the shallowest feed that identifies the event cleanly. Order entry (new, cancel, replace, reject and execution-report messages), public market data, drop copies (private execution and routing records) and clearing feeds are separate message streams.

Velu, Hardy and Nehren describe the same path: client orders reach a broker in FIX, an order management system keeps each order's state, an outbound gateway translates child orders into the venues' protocols (``Essentially all of them support a FIX protocol, but in most cases they also support a much faster `native' protocol'') and captures acknowledgements and executions, and clearing and settlement follow \citep[chap.~11, sec.~11.1, pp.~383--386]{s3-velu}. Drop copies and clearing-feed contents are an additional interpretation, not in the source; the FIX order-state and execution-report values are in \hyperref[mech:order-lifecycle]{the order lifecycle}.

\qffield{Pitfalls} The price field is divided by 10000, not 1000 (p.~37). Reconstructing the full book requires consolidation, and reception timestamps include latency. The Table~1.10 layout is the book's illustration; a real feed's specification, which changes, governs.

\qffield{See also} \hyperref[data:trading-reference-data]{trading reference data}; \hyperref[def:order-exposure-decisions]{order exposure decisions} (hidden and reserve orders); \hyperref[mod:static-limit-order-depth]{static limit-order depth}; \hyperref[meth:information-driven-bars]{information-driven bars}.

\qffield{Further reading} \citep{wide-holden-jacobsen}, published version and Internet Appendix (MTAQ program).

\end{qfdata}

\begin{qfdata}[Derived market statistics: daily activity data, intraday microstructure statistics, pre-computed aggregates and binned data]\label{data:derived-market-statistics}
\qffield{Statement} Velu, Hardy and Nehren group the statistics most strategies derive from binned or daily data \citep[chap.~1, sec.~1.4.3, pp.~39--42]{s3-velu}. \emph{Daily activity}: open, high, low, close and previous close (adjusted for corporate actions), the last trade before the close, volume with its lit and dark split, auction volume and price, VWAP, short interest, days-to-cover and borrow utilisation, futures open interest and basis, index-level data, options open interest, Greeks and implied volatility, and cross-asset rates, CDS spreads and the dollar index (pp.~39--40).

\emph{Intraday microstructure}: number and frequency of trades and quote updates, top-of-book size, depth, spread, trade size and ticking time, whose daily distributions can start an algorithm's day and be updated ``through online Bayesian updates'' (p.~41). \emph{Pre-computed aggregates}: $X$-day ADV and auction volume, volatility, beta (plain or asymmetric) and correlation matrices (p.~41). \emph{Binned data} discretise the day into bins of a few seconds to 30 minutes with the same features; volume and spread profiles ``are rarely built with a granularity finer than one minute'' to avoid noise from frictions (p.~42).

\qffield{Conditions} Practitioner descriptions as of about 2018--2020. Aggregates ``need to fully support the peculiarities described in the Reference Data section'': special days skew distributions, and asynchronous markets distort beta and correlation (p.~41). Binned data allow discrete time-series methods and reduce dimension (390 one-minute rows per U.S. stock-day against millions of messages), but relationships stable on daily closes become noisier at higher frequency, and illiquid names leave empty bins (p.~42). Interpretation: each derived field must carry the time it became available, for example short interest only after its publication date, and pre-computed $X$-day aggregates must be computed through the decision date, not the end of the sample.

\emph{Short-sale availability.} Nasdaq's daily Regulation SHO threshold-security files are one public point-in-time input. The files examined, for 10 January 2005, 2 January 2008 and 15 October 2025, are each a pipe-delimited list with a threshold flag and a trailing creation stamp (\texttt{20050110231003}, created after the close), and the 2025 file adds a Rule~3210 flag \citep[files \texttt{nasdaqth20050110}, \texttt{nasdaqth20080102} and \texttt{nasdaqth20251015}]{nasdaq-regsho-threshold}. A threshold security is one with substantial fails to deliver, which clearing participants must close out under Rule 203(b)(3) of Regulation SHO; the rule's compliance date was 3 January 2005 \citep[Discussion]{finra-ntm-04-93}. The list measures delivery failures, not borrow cost or availability. The sources here give no public point-in-time record of hard-to-borrow status or borrow fees, so those remain the gap Narang names.

\qffield{Practice} Velu, Hardy and Nehren's uses are judgements to test, not results: the high--low range ``usually gives a better indication of market sentiment than the simple close-to-close return''; ``a large short interest usually indicates a bearish view from institutional investors''; days-to-cover sizes a potential short squeeze; futures basis against one's dividend estimates may show mispricing (pp.~39--40).

\qffield{Pitfalls} Daily activity summaries differ from binned data (p.~42); both require consistent reference data when aggregating.

\qffield{See also} \hyperref[meth:garman-klass-volatility]{Garman--Klass volatility} (range-based estimators); \hyperref[def:spread-measures]{spread measures}; \hyperref[def:securities-lending-fee-rebate]{securities lending fees and rebates} (utilisation, hard to borrow); \hyperref[meth:activity-sampled-bars]{activity-sampled bars}; \hyperref[data:trading-reference-data]{trading reference data}.

\end{qfdata}

\begin{qfdata}[Holdings and insider filings: Form 13F, Schedules 13D and 13G, Forms 3 and 4, their deadlines, and when their contents can enter a decision-time panel]\label{data:holdings-insider-filing-lags}
\qffield{Statement} U.S. deadlines in force on 15 October 2026 (eCFR point-in-time text):
\begin{itemize}
\item \emph{Form 13F}: an institutional investment manager exercising investment discretion over section 13(f) securities with an aggregate fair market value of at least \$100{,}000{,}000 ``on the last trading day of any month of any calendar year'' files ``within 45 days after the last day of such calendar year and within 45 days after the last day of each of the first three calendar quarters of the subsequent calendar year'' \citep[17 CFR 240.13f-1(a)(1)]{ecfr-17cfr240-filing-rules}; only securities on the Commission's 13(f) list count and are reported (para.~(c)).
\item \emph{Schedule 13D}: a person who becomes the beneficial owner of more than five percent of a covered class files ``within five business days after the date of the acquisition'' (240.13d-1(a)); a material change, including an acquisition or disposition of one percent or more of the class, requires an amendment ``within two business days after the date of such change'' (240.13d-2(a)).
\item \emph{Schedule 13G}: qualified institutions file within 45 days after the quarter in which they become obligated, or within five business days after the end of the first month in which ownership exceeds 10 percent (240.13d-1(b)(2)); passive investors under 20 percent within five business days after the acquisition (240.13d-1(c)).
\item \emph{Forms 3 and 4}: insiders (directors, officers and more-than-10-percent owners) file Form~3 ``within 10 days after he or she becomes such beneficial owner, director, or officer'' and report a change ``before the end of the second business day following the day on which the subject transaction has been executed'' \citep[15 U.S.C. 78p(a)(2)(B)--(C), 2023 edition]{usc-15-78p}, on Form~4 (240.16a-3(a) and (g)(1)).
\end{itemize}
Under Rule 0-3(a), papers whose last timely day ``falls on a Saturday, Sunday or holiday \dots\ may be filed on the first business day following''; Regulation S-T sets the \emph{filing date} of an electronic submission. Electronic filings may be transmitted on business days from 6~a.m.\ to 10~p.m.\ Eastern \citep[17 CFR 232.12(c)]{ecfr-17cfr232-filing-date}. A direct transmission ``commencing on or before 5:30 p.m.'' Eastern is ``deemed filed on the same business day'', and one commencing later is deemed filed the next business day (232.13(a)(2)). Forms 3, 4, 5 and 144 and Schedules 13D, 13G and 14N, with their amendments, are deemed filed the same day up to 10~p.m.\ (232.13(a)(4)). The filing date can be adjusted for technical difficulties (232.13(b)).

\qffield{Conditions} \emph{History.} The 13D and 13G deadlines above come from the Commission's 2023 amendments, effective 5 February 2024, which shortened the initial 13D deadline from ``within 10 days after the acquisition'' to five business days; compliance with the revised 13G deadlines was not required before 30 September 2024 \citep[88~FR 76896, 7 November 2023, pp.~76897 and~76942]{fr-2023-22678}. A panel covering earlier years must apply the deadline in force at each filing date.

\emph{Deemed execution dates.} For trades under a plan or instruction that the insider does not time, the date the broker or plan administrator notifies the insider counts as the execution date, but no later than the third business day after the trade (240.16a-3(g)(2)--(4)); such Form~4 filings can therefore come up to five business days after the trade. Velu, Hardy and Nehren's footnotes (``within 10 days'' for 13D; Form~4 ``within 2 days''; 13F managers ``with over \$100 million in assets'' reporting quarterly) \citep[chap.~1, sec.~1.4, p.~43, nn.~31--32]{s3-velu} predate the 2024 change and simplify the 13F test, which counts 13(f) securities, not assets.

\qffield{Practice} \emph{Four timestamps, not one.} A decision-time panel should keep these apart for each filing:
\begin{itemize}
\item EDGAR's \emph{acceptance} time of the submission;
\item the official \emph{filing date} under Rule 232.13. For a 10-K or 13F transmitted after 5:30~p.m.\ it is the next business day, which can make a deadline-day 13F late; a Form~4 is still filed the same day up to 10~p.m.;
\item the time of \emph{public dissemination} on EDGAR;
\item the time the user's own system \emph{first observed} the filing.
\end{itemize}
EDGAR's page \emph{Determine the Status of My Filing} (dated 4 June 2024, last reviewed 31 October 2024) states that most submissions that begin transmission after 5:30~p.m.\ ``will receive a filing date of 6:00 a.m. ET the next business day and will not be disseminated by EDGAR until the next business day'', while Forms 3, 4, 5 and 144 (with amendments) and some MEF and POS~462B forms transmitted before 10~p.m.\ get that day's filing date and ``will be disseminated by EDGAR until 10:00 p.m. ET on that date'' \citep[sec.\ \emph{Filing Date of Electronically-transmitted Submissions}]{sec-edgar-filing-status}. The page gives a day, not a time, for dissemination, and no source bounds it from above.

So the start of the next business day is not a conservative availability time: an after-cutoff 10-K or 13F cannot be public at 00:00, and may not be public until some later, unstated hour. The panel keeps two quantities. \emph{Availability} is the recorded dissemination time, else the recorded first-observed time; with neither, it is \emph{unknown} and the filing does not enter a decision-time panel on that record. A \emph{lower bound} (acceptance; for an after-cutoff submission not on the same-day list, 6:00~a.m.\ on the next business day, after any holiday) only detects certain look-ahead: use before it is a leak, use after it is not thereby safe, and a recorded first observation earlier than the bound is a data error.

Schedules 13D, 13G and 14N are deemed filed the same day up to 10~p.m.\ under 232.13(a)(4), but the page's same-day dissemination list omits them, so their lower bound is acceptance alone. Any fixed delay a user adopts instead (a decision-time convention) is an assumption to state and vary, not a rule. The deadline gives only the latest date, and for 13F the reported positions can be up to about 45 days (plus a weekend or holiday roll) old when public.

Worked example: positions at 30 September 2026 are due by 14 November 2026, a Saturday, so by Monday 16 November; a Form~4 for a trade on Friday 20 November 2026 is due by Tuesday 24 November; a 13D for an acquisition on Monday 23 November 2026 is due by Tuesday 1 December, Thanksgiving (26 November) not counting. Amendments (13F/A, 13D/A) restate earlier filings and keep their own four timestamps; they never overwrite the original's (\hyperref[data:trading-reference-data]{trading reference data}). The same holds for information first withheld and later made public. Section~13(f)(4) lets the Commission ``delay or prevent public disclosure'' of 13F information \citep[15 U.S.C. 78m(f)(4), 2023 edition]{usc-15-78m}, and confidential-treatment requests follow Rule 24b-2, which requires them in paper \citep[17 CFR 240.24b-2, preamble]{ecfr-17cfr240-24b-2}.

Positions released later are known only from their release, so they enter the panel at the later filing's timestamps, not at the original period. These filing-date rules and timestamps are checked on synthetic records.

\qffield{Pitfalls} The entry covers deadlines, not the data sets: confidential treatment of 13F positions, the structure of the SEC's 13F and N-PORT data sets and the forms' instructions are not covered. EDGAR's dissemination timing is stated here only from the status page above, which gives days, not times; the EDGAR Filer Manual is not covered. The 10~p.m.\ cut-off for 13D and 13G that the 2023 release mentions (88~FR 76897) is now confirmed in 232.13(a)(4). Federal holidays in the example are those of 5 U.S.C. 6103.

\qffield{See also} \hyperref[data:trading-reference-data]{trading reference data}; \hyperref[def:inside-information]{inside information} (Harris's account of the insider-reporting rules as of 2002); \hyperref[def:financial-data-types]{four types of financial data}.

\qffield{Further reading} The SEC's Form 13F frequently asked questions and the 13F list of section 13(f) securities.

\end{qfdata}

\begin{qfmethod}[Bitemporal reference tables: actual (valid) time against record (transaction) time for identifiers, corporate actions and vendor corrections]\label{meth:bitemporal-reference-tables}
\qffield{Statement} Fowler separates two time dimensions of a history \citep[sec.~``The Two Dimensions'']{s8-bitemporal}:
\begin{itemize}
\item \emph{actual} history, ``what history should be given perfect transmission of information'';
\item \emph{record} history, ``how our knowledge of history changes''.
\end{itemize}
He notes that these are also called \emph{valid} (or effective) and \emph{transaction} time, Snodgrass's terms used in SQL:2011 (n.~1). A value is then read as $\text{value}(\text{actual date},\text{record date})$, such as ``on Mar 25th, we thought Sally's salary on Feb 25th was \$6500'' (sec.~``The Two Dimensions''). The other properties he gives:
\begin{itemize}
\item Retroactive changes rewrite actual history, but ``record history itself is append only'' (sec.~``Append-only History'').
\item Each agent (department, vendor) has its own record times (sec.~``Perspectives for Record Time'').
\item Storage can be a table with record and actual start and end dates, or \emph{event sourcing}. Events must then carry both an actual date and a record date (sec.~``Storing and Processing Bitemporal Histories'').
\end{itemize}

\qffield{Notation} Actual time $a$ (when a fact held in the world); record time $r$ (when this system or vendor recorded it); a query is $(a,r)$; ``as known at decision time $t$'' sets $r=t$.

\qffield{Conditions} Concepts as published 2021. \emph{Applied to trading reference data} (Interpretation, not in Fowler). Several kinds of data need both axes:
\begin{itemize}
\item identifier changes, such as tickers and CUSIP or ISIN reassignments;
\item corporate actions and their adjustment factors;
\item index membership;
\item vendor corrections of prices and fundamentals;
\item filings and their amendments, whose four timestamps are in \hyperref[data:holdings-insider-filing-lags]{holdings and insider filings}.
\end{itemize}
A backtest at decision time $t$ must query $r\le t$. Today's corrected values use $r=$ now and leak later knowledge into the past. Fowler's caution applies too. Bitemporality ``does complicate a system quite significantly'' and is avoidable when no action depends on retroactively changed values; but a backtest's actions do depend on what was known (sec.~``Using Bitemporality'').

\qffield{Practice} Synthetic check. A vendor records a 2-for-1 split effective 2 March 2026 on 3 March and corrects it to 3-for-2 on 20 March. Querying the factor for 5 March:
\begin{itemize}
\item as known on 10 March, it returns 2.0, what a backtest dated then saw;
\item as known now, it returns 1.5;
\item as known on 2 March, it returns 1.0;
\item the correction closes record ranges and appends rows, so earlier rows survive unchanged.
\end{itemize}
The check verifies this illustration's logic only; vendors' actual correction histories were not examined.

\qffield{Pitfalls} Many vendor data sets keep only the latest actual history, and some keep no record history, so a recorded ``as of'' date can be a later reconstruction. Fowler's ``actual time is always at or before record time'' is stated for history; recorded future events, such as an announced split, reverse it (n.~2).

\qffield{See also} \hyperref[data:trading-reference-data]{trading reference data}; \hyperref[data:holdings-insider-filing-lags]{holdings and insider filings}; \hyperref[pit:label-weight-leakage]{time-validity of labels and weights}; \hyperref[meth:time-valid-ml-experiment-record]{time-valid experiment records}.

\qffield{Further reading} Snodgrass's \emph{Developing Time-Oriented Database Applications in SQL}, which Fowler cites.

\end{qfmethod}
